\chapter{Epochen}

\startframedtext[width=\dimexpr\textwidth+\rightmarginwidth+\rightmargindistance\relax,frame=off,offset=none] % This spans textwidth + margin width
\startplacetable[location={here},title={Einordnung}]
\setupTABLE[c][1][background=color,backgroundcolor=gray,style=\bf]
\setupTABLE[r][1][background=color,backgroundcolor=gray,style=\bf]
\bTABLE[option=stretch,split=yes]
\startCSV
, "1945                                                                                       ", "1968                                                                   ", "1989                                                            "
"Literarische Strömungen", "Beginn der Hinwendung zur Postmoderne, Trümmerliteratur als Vorläufer              ", "Intensivierung postmoderner Tendenzen, neue Subjektivität, Popliteratur", "Vollendung der Postmoderne, Wendeliteratur, Post-DDR-Literatur  "
"Sprachliche Merkmale   ", "Beginn der Skepsis gegenüber “großen Erzählungen”, Hinwendung zu fragmentierter Erzählweise", "Experimentieren mit Sprache, Dekonstruktion von Sinn, Ironie           ", "Intertextualität, Pluralität der Stimmen, Metafiktion           "
"Zentrale Themen        ", "Auseinandersetzung mit Schuld, Identitätskrise, Suche nach neuer Bedeutung                 ", "Gesellschaftskritik, Entlarvung von Ideologien, Subversion             ", "Reflexion über Geschichte, Identität, Auflösung von Gewissheiten"
\stopCSV
\eTABLE
\stopplacetable
\stopframedtext

\startplacetable[location={here,split},title={Zusammenfassung der Epochen \cite[num][Epochen_stroemung_deutschabi_truemmer]}]
\setupTABLE[c][1][background=color,backgroundcolor=gray,style=\bf]
\setupTABLE[each][offset=1mm]
\bTABLE[split=yes,option=stretch,textwidth=\dimexpr\textwidth+\rightmarginwidth+\rightmargindistance\relax]
\bTR
\bTD[nc=2] Nachkriegsliteratur / Trümmerliteratur \eTD
\eTR
\bTR \bTD Ab Seite: \eTD \bTD \at[section:epochen:nachkriegsliteratur] \eTD \eTR
\bTR
\bTD Authoren: \eTD \bTD Wolfgang Borchert (\at{s. S.}[personen:wolfgang_borchert]) \eTD
\eTR
\bTR
\bTD Zeit \eTD \bTD 1945 - 1965 \eTD
\eTR
\bTR
\bTD Hintergrund \eTD \bTD Zweiter Weltkrieg \eTD
\eTR
\bTR
	\bTD
	Motive, Themen und Ideen 
	\eTD

	\bTD
	Verarbeitung der Geschichte sowie der „Trümmerlandschaft“, Ablehnung von Vergessen / Verdrängung des Kriegsgeschehens, Wiedergabe der Wirklichkeit, Ablehnung von NS-Sprache und derer Prägungen, Literaturströmung in der Bevölkerung / Gesellschaft unbeliebt, zum Teil auch Verarbeitung von Kriegserlebnissen und -ereignissen.
	\eTD
\eTR
\bTR
	\bTD
	Textgattungen, -sorten und formale Merkmale
	\eTD
	\bTD
	Einfache und besonders deutliche Sprache um Verwechslungen oder Anspielungen mit NS-Begriffen o.ä. zu vermeiden
	\eTD
\eTR

\bTR
\bTD[nc=2] Postmoderne \eTD
\eTR
\bTR \bTD Ab Seite: \eTD \bTD \at[section:epochen:nachkriegsliteratur] \eTD \eTR
\bTR
\bTD Authoren: \eTD \bTD Wolfgang Borchert (\at{s. S.}[personen:wolfgang_borchert]) \eTD
\eTR
\bTR
\bTD Zeit \eTD \bTD 1945 - 1965 \eTD
\eTR
\bTR
\bTD Hintergrund \eTD \bTD Zweiter Weltkrieg \eTD
\eTR
\bTR
	\bTD
	Motive, Themen und Ideen 
	\eTD

	\bTD
	Verarbeitung der Geschichte sowie der „Trümmerlandschaft“, Ablehnung von Vergessen / Verdrängung des Kriegsgeschehens, Wiedergabe der Wirklichkeit, Ablehnung von NS-Sprache und derer Prägungen, Literaturströmung in der Bevölkerung / Gesellschaft unbeliebt, zum Teil auch Verarbeitung von Kriegserlebnissen und -ereignissen.
	\eTD
\eTR
\bTR
	\bTD
	Textgattungen, -sorten und formale Merkmale
	\eTD
	\bTD
	Einfache und besonders deutliche Sprache um Verwechslungen oder Anspielungen mit NS-Begriffen o.ä. zu vermeiden
	\eTD
\eTR

\bTR
\bTD[nc=2] Nachkriegsliteratur / Trümmerliteratur \eTD
\eTR
\bTR \bTD Ab Seite: \eTD \bTD \at[section:epochen:nachkriegsliteratur] \eTD \eTR
\bTR
\bTD Authoren: \eTD \bTD Wolfgang Borchert (\at{s. S.}[personen:wolfgang_borchert]) \eTD
\eTR
\bTR
\bTD Zeit \eTD \bTD 1945 - 1965 \eTD
\eTR
\bTR
\bTD Hintergrund \eTD \bTD Zweiter Weltkrieg \eTD
\eTR
\bTR
	\bTD
	Motive, Themen und Ideen 
	\eTD

	\bTD
	Verarbeitung der Geschichte sowie der „Trümmerlandschaft“, Ablehnung von Vergessen / Verdrängung des Kriegsgeschehens, Wiedergabe der Wirklichkeit, Ablehnung von NS-Sprache und derer Prägungen, Literaturströmung in der Bevölkerung / Gesellschaft unbeliebt, zum Teil auch Verarbeitung von Kriegserlebnissen und -ereignissen.
	\eTD
\eTR
\bTR
	\bTD
	Textgattungen, -sorten und formale Merkmale
	\eTD
	\bTD
	Einfache und besonders deutliche Sprache um Verwechslungen oder Anspielungen mit NS-Begriffen o.ä. zu vermeiden
	\eTD
\eTR

\bTR
\bTD[nc=2] Nachkriegsliteratur / Trümmerliteratur \eTD
\eTR
\bTR \bTD Ab Seite: \eTD \bTD \at[section:epochen:nachkriegsliteratur] \eTD \eTR
\bTR
\bTD Authoren: \eTD \bTD Wolfgang Borchert (\at{s. S.}[personen:wolfgang_borchert]) \eTD
\eTR
\bTR
\bTD Zeit \eTD \bTD 1945 - 1965 \eTD
\eTR
\bTR
\bTD Hintergrund \eTD \bTD Zweiter Weltkrieg \eTD
\eTR
\bTR
	\bTD
	Motive, Themen und Ideen 
	\eTD

	\bTD
	Verarbeitung der Geschichte sowie der „Trümmerlandschaft“, Ablehnung von Vergessen / Verdrängung des Kriegsgeschehens, Wiedergabe der Wirklichkeit, Ablehnung von NS-Sprache und derer Prägungen, Literaturströmung in der Bevölkerung / Gesellschaft unbeliebt, zum Teil auch Verarbeitung von Kriegserlebnissen und -ereignissen.
	\eTD
\eTR
\bTR
	\bTD
	Textgattungen, -sorten und formale Merkmale
	\eTD
	\bTD
	Einfache und besonders deutliche Sprache um Verwechslungen oder Anspielungen mit NS-Begriffen o.ä. zu vermeiden
	\eTD
\eTR
\eTABLE
\stopplacetable
% \stopframedtext

\section[section:epochen:nachkriegsliteratur]{Nachkriegsliteratur}

Umfasst drei Schlagwörter \cite[authornum][boell2007werke]:
\startitemize[packed]
\item Krigesliteratur
\item Heimkehrliteratur
\item Trümmerliteratur
\stopitemize

Man wollte nicht die Augen vor der Grausammkeit verschließen und "Blinde Kuh spielen" \cite[authornum][boell2007werke]

% \startitemize[packed]
% \item \cite[default][boell2007werke]
% \item \cite[category][boell2007werke]
% \item \cite[entry][boell2007werke]
% \item \cite[short][boell2007werke]
% \item \cite[page][boell2007werke]
% \item \cite[num][boell2007werke]
% \item \cite[textnum][boell2007werke]
% \item \cite[year][boell2007werke]
% \item \cite[index][boell2007werke]
% \item \cite[tag][boell2007werke]
% \item \cite[keywords][boell2007werke]
% \item \cite[author][boell2007werke]
% \item \cite[authoryears][boell2007werke]
% \item \cite[authornum][boell2007werke]
% \item \cite[authoryear][boell2007werke]
% \stopitemize

\subsection{Historischer Hintergrund - 2. Weltkrieg}

\subsection{Motive, Themen und Ideen}
